\chapter{Đánh giá hệ thống}
\label{chap:danhgia}

Chương này trình bày phương pháp và kết quả đánh giá Farmery theo sáu trụ cột: kiểm thử chức năng, chất lượng mã nguồn, hiệu năng và chịu tải phía máy chủ, an toàn bảo mật, khả dụng giao diện và chất lượng mô-đun AI. Sáu trụ cột gộp vào bốn nhóm mục tiêu ở mục~\ref{sec:muctieudanhgia}: chức năng và chất lượng mã thuộc nhóm đánh giá chức năng; hiệu năng và bảo mật thuộc nhóm hiệu năng và độ tin cậy.

\section{Mục tiêu và phương pháp đánh giá}
\label{sec:muctieuphuongphap}

Mọi ngưỡng đối chiếu được xác định \emph{trước} khi hiện thực --- Bảng~\ref{tab:nguongNFR} cho yêu cầu phi chức năng và Bảng~\ref{tab:nguongAI} cho mô-đun AI --- để không chọn chỉ số sau khi đã biết kết quả.
\begin{itemize}
    \item \textbf{Chất lượng mã nguồn:} typecheck, ESLint và độ phủ Jest cho phần nghiệp vụ lõi, chạy tự động qua GitHub Actions với cổng độ phủ chặn mã không đạt ngưỡng (mục~\ref{sec:moitruong}).
    \item \textbf{Chức năng:} kiểm thử hộp đen các nhóm yêu cầu có phần hiện thực (FR9 chỉ thiết kế), chỉ kiểm phần không mang nhãn chỉ thiết kế, ưu tiên luồng đi xuyên nhiều phân hệ: hợp đồng sỉ và đơn thu mua cùng giành hàng một lô, thu mua vào kho.
    \item \textbf{Hiệu năng và chống bán vượt:} k6 (kịch bản JavaScript, nhẹ hơn JMeter, gắn được vào tích hợp liên tục) đo p95, thông lượng, và giả lập nhiều yêu cầu đồng thời tranh một lô có tồn giới hạn.
    \item \textbf{Bảo mật:} kiểm thử tự động ma trận RBAC trên mọi endpoint nghiệp vụ; rà soát dữ liệu cá nhân theo Luật Bảo vệ dữ liệu cá nhân 2025 (NFR8.2).
    \item \textbf{Khả dụng giao diện:} giao diện đáp ứng trên ba kích thước màn hình; ba chỉ số Core Web Vitals bằng Lighthouse; buổi thử ba tác vụ với nhóm nhà cung cấp đo bốn chỉ số hành vi; axe-core; danh sách kiểm tra PWA trên điện thoại thật.
    \item \textbf{Mô-đun AI:} đánh giá offline bằng chỉ số riêng từng bài toán; dự báo giá đánh giá trên dữ liệu giá công khai chứ không phải dữ liệu giao dịch của nền tảng.
\end{itemize}

% TODO: nêu điều kiện đo của toàn chương (cấu hình máy, phiên bản công cụ, kích thước tập dữ liệu mô phỏng, ngày thực hiện) để kết quả tái lập được.

\section{Đánh giá chức năng}
\label{sec:danhgiachucnang}

Kiểm thử hộp đen theo bộ test case bám các nhóm yêu cầu, với ba kịch bản trọng tâm: thu mua từ chào giá lô tới nhập kho và mở bán; bán sỉ từ yêu cầu báo giá, một vòng phản hồi giá tới giao một đợt; bán lẻ từ đặt hàng tới đổi trả.

% TODO: bảng tổng hợp đạt hoặc không đạt theo từng nhóm yêu cầu chức năng; bảng test case chi tiết ở mục~\ref{sec:kichbankiemthu}.

\section{Đánh giá phía máy chủ (back-end)}
\label{sec:danhgiahieunang}

\noindent\textbf{Hiệu năng và chịu tải.} Kịch bản tải tăng dần số người dùng đồng thời thực hiện chuỗi: tìm kiếm, xem chi tiết và nhật ký canh tác, thêm giỏ, tạo đơn. Theo Bảng~\ref{tab:nguongNFR}, báo cáo p95 dưới 2 giây ở tải thông thường, và ở mức mục tiêu 100 người dùng đồng thời thì thời gian phản hồi trung vị dưới 3 giây cùng tỷ lệ lỗi dưới 1\% (NFR1.1). Song song, bật \texttt{pg\_stat\_statements} để tìm truy vấn chậm nhất và truy vấn bị gọi lặp.

\noindent\textbf{Chống bán vượt tồn kho.} Để kiểm chứng NFR1.2 và NFR4.1, kịch bản tranh chấp khởi tạo dòng tồn có số lượng giới hạn rồi kích hoạt đồng thời nhiều yêu cầu vượt số đó, theo hai kịch bản: nhiều đơn lẻ cùng đặt trên dòng tồn của Farmery; hợp đồng sỉ và đơn thu mua cùng giành dòng tồn của nhà cung cấp. Cần chứng minh khóa bi quan (\texttt{SELECT \ldots{} FOR UPDATE}) chặn mọi trường hợp bán vượt ở cả hai kịch bản.

\noindent\textbf{Phân quyền và bảo mật.} Script gọi endpoint nghiệp vụ bằng token của vai trò không đủ quyền (ví dụ \texttt{buyer\_retail} gọi endpoint đăng bán, token một vai trò nội bộ gọi endpoint vai trò khác) và xác nhận mọi lời gọi bị chặn với mã 403 (NFR3.1); rà soát cơ sở dữ liệu để xác nhận mã hóa dữ liệu định danh và băm mật khẩu (NFR3.2, NFR3.3).

\noindent\textbf{Độ tin cậy và các yêu cầu còn lại.} Thử khôi phục từ bản sao lưu ít nhất một lần trong giai đoạn đánh giá (NFR4.2). Các yêu cầu định tính hoặc chỉ thiết kế --- NFR2.1, NFR2.2, NFR3.4--NFR3.6, NFR5.1, NFR5.2, NFR7.1, NFR8.1--NFR8.5, NFR9.1 --- được nghiệm thu bằng danh sách kiểm tra hoặc rà soát thiết kế chứ không đo bằng chỉ số, đúng như mục~\ref{sec:yeucauphichucnang}.

% TODO: điền số đo từng mức tải kèm kết luận so với NFR1.1 và năm truy vấn tốn thời gian nhất; số yêu cầu đồng thời, số lượng ban đầu, số được chấp nhận và bị từ chối, tỷ lệ bán vượt thực đo; số endpoint đã kiểm, tỷ lệ chặn đúng, kết quả rà soát mã hóa.

\section{Đánh giá phía giao diện (front-end)}
\label{sec:danhgiaux}

\noindent\textbf{Tính đáp ứng.} Kiểm thử trên điện thoại, máy tính bảng, máy tính để bàn với hai trình duyệt phổ biến ở Việt Nam, các màn hình chính: trang chủ, chi tiết sản phẩm, tra cứu QR, giỏ hàng, trang quản lý của nhà cung cấp; tiêu chí không tràn ngang, không che nút thao tác (NFR7.2). Lighthouse đo Core Web Vitals với ngưỡng ``tốt'': LCP không quá 2,5 giây, INP không quá 200 ms, CLS không quá 0,1 \parencite{webdev_vitals}. PWA không đo bằng Lighthouse (đã bỏ hạng mục này từ phiên bản 12) mà kiểm theo danh sách ở Phụ lục~A, mục A.6 (NFR7.3).

\noindent\textbf{Khả dụng với nhóm nhà cung cấp.} Nhóm đại diện gồm nông dân và thành viên ban điều hành hợp tác xã thử ba tác vụ khó và thường xuyên nhất: đăng một lô hàng kèm ảnh, ghi nhật ký canh tác cho lô, tìm và mở QR của lô. Bốn chỉ số hành vi: tỷ lệ hoàn thành không cần hỗ trợ (ngưỡng NFR6.1--NFR6.3), thời gian hoàn thành, số lần nhờ giúp, số lỗi phải làm lại; không dùng bảng câu hỏi chuẩn hóa vì điểm cảm nhận không chỉ ra chỗ người dùng vướng. Bổ sung axe-core theo WCAG mức AA, có nêu giới hạn của kiểm tra tự động. Kịch bản, phiếu ghi và cách tính ở Phụ lục~A.

% TODO: điền kết quả từng nhóm thiết bị, ba chỉ số Core Web Vitals, kết quả danh sách PWA; số người tham gia, ngày, địa điểm, bốn chỉ số từng tác vụ, chỗ vướng, số lỗi axe-core trước và sau khi sửa.

\section{Đánh giá mô-đun AI}
\label{sec:danhgiaAI}

Nghiệm thu theo chỉ số ở Bảng~\ref{tab:nguongAI}: \emph{tìm kiếm ngữ nghĩa} đo nDCG@10, Recall@10 trên bộ truy vấn gán nhãn tay so với tìm từ khóa làm đường cơ sở, kèm p95; \emph{sản phẩm tương tự} đo Precision@5 trên tập sản phẩm gán nhãn tay (CTR qua A/B chỉ áp dụng khi bật cá nhân hóa, ngoài phần hiện thực); \emph{dự báo giá} đo MAPE trên tập kiểm tra tách theo thời gian, dùng chuỗi giá theo tháng của FAOSTAT (đến 01/2025, đã nối tại điểm gãy mức giá 01/2022) cho bốn mặt hàng khớp sản phẩm của nền tảng là khoai tây, nấm, chè, cà phê --- dâu tây, cà chua và các mặt hàng đặc trưng Đà Lạt không có chuỗi đủ dài nên không đo được; mô hình chạy offline và chỉ đánh giá bằng MAPE. Hai mô hình đối chứng ở Phụ lục~F --- Multilingual-E5 cho mô hình nhúng và gradient boosting cho dự báo giá --- được đo cùng chỉ số trên cùng tập kiểm tra. Giám sát bất thường dùng luật ngưỡng (Chương~\ref{chap:coso}) nên được nghiệm thu bằng việc kiểm tra luật kích hoạt đúng trên tập tình huống dựng sẵn.

% TODO: điền kết quả thực nghiệm, kết luận đạt hoặc không đạt, kích thước và nguồn từng tập dữ liệu. Cuối chương tổng hợp ưu điểm và hạn chế, nêu cả chỉ số chưa đạt ngưỡng kèm nguyên nhân (hội đồng quan tâm phần này hơn danh sách chỉ số đã đạt) và phản ánh lại vào mục Hạn chế ở Chương~\ref{chap:ketluan}.
